\documentclass[10pt,a4paper]{amsart}
\usepackage[margin=19mm]{geometry}
\usepackage{amsmath,amssymb,mathtools}
\usepackage[scr=rsfs,cal=euler]{mathalfa}
\usepackage{xfrac}
\usepackage[unicode,hidelinks,bookmarks=false]{hyperref}
\newcommand{\ie}{i.e.\ }
\newcommand{\cf}{cf.\ }
\newcommand{\resp}{respectively\ }
\newcommand{\Subsection}{\subsection}
\providecommand{\nobreakdash}{\nobreak\hbox{-}}
\setlength{\emergencystretch}{2em}
\multlinegap0pt
\usepackage{mathtools}
\usepackage{tikz}
\usepackage{amssymb}
\usepackage{bm}
\usepackage{dsfont}
\usepackage{enumerate}
\usepackage[shortlabels]{enumitem}
\newtheorem{lemma}{Lemma}
\newcommand{\CC}{\textnormal{C}}
\newcommand{\mA}{\mathcal{A}}
\newcommand{\mC}{\mathcal{C}}
\newcommand{\mU}{\mathcal{U}}
\title{\textbf{Capacity of Non-Separable Networks with Restricted Adversaries}}
\author{Christopher Hojny, Altan B. K\i l\i\c{c}, Sascha Kurz, Alberto Ravagnani}
\date{Source: arXiv:2603.08246v1 (CC BY 4.0)}
\begin{document}
\maketitle

\noindent\emph{Source and licence.} \textbf{Capacity of Non-Separable Networks with Restricted Adversaries}. Version arXiv:2603.08246v1 (CC BY 4.0), distributed by arXiv under the Creative Commons Attribution 4.0 International licence, http://creativecommons.org/licenses/by/4.0/, as recorded in the arXiv licence metadata for this exact version and shown on the abstract page https://arxiv.org/abs/2603.08246v1. Full licence notice: \texttt{licenses/arxiv-2603.08246-CC-BY-4.0.txt}.\par\smallskip

\noindent\textit{Editorial scope.} Counted unit: the article's lemma lemma1 (source0.tex lines 703--706). The article's complete original proof of it is delivered with it (source0.tex lines 707--730, one \texttt{proof} environment of 24 lines). No mathematics is written, repaired or re-derived: the statement and the proof are the article's own lines, in the article's own notation, and no retained line is substituted or re-flowed.  No further source-local context is needed: every cross-reference of every retained line resolves inside the two delivered spans.  Nothing was omitted from any retained span, so the package carries no omission record.  No retained line contains a citation, so the wrapper carries no bibliography and no bibliography entry.  No retained line contains a figure environment, an image-inclusion command or drawing code, so no image is delivered and the package declares no asset dependency.  Editorial formatting only, in addition: an AMS \texttt{amsart} preamble that restates the article's own declarations and macros used by the retained lines, namely the environments \texttt{lemma} on separate counters, together with the standard packages those lines need.  The retained environments print in this document as Lemma 1 in order of appearance, whereas the article prints the same items with the numbering of its own full document.  The complete native source of the article as distributed in the arXiv e-print for this exact version is bundled unchanged as \texttt{sources/195991/source0.tex}.
\begin{lemma}
\label{lemma1}
Let $m$ be a nonnegative integer and $t=2m+1$. We have $$\CC_1(\mathfrak{E}_t,\mA,\mU_S,t) \ge \log_{|\mA|}\left(\left\lfloor\tfrac{|\mathcal{A}|-1}{m+1}\right\rfloor\right).$$
\end{lemma}

\begin{proof}
Let $\mathcal{A}'\subseteq\mathcal{A}$ be an arbitrary subset of the alphabet of cardinality $a'\coloneqq\left\lfloor(|\mathcal{A}|-1)/(m+1)\right\rfloor$.
As outer code, we choose a $(2t+1)$-fold repetition code over the sub-alphabet $\mathcal{A}'$, i.e., $$\mC\coloneqq\left\{(c,\dots,c)\in\mathcal{A}^{2t+1}\,\mid\,c\in\mathcal{A}'\right\},$$ of cardinality $a'$.


For the function of $V_2$, we inject
the set of output states of $V_2$ into $\mathcal{A}$. There are $2m+2$ incoming edges of $V_2$. To have a unique majority winner that occurs more than half of those edges, call it $x$, the symbol $x$ must occur on at least $m+2$ of those edges. We have~$(2m+2)-(m+2)+1=m+1$. Thus, the injection we are looking for is of the form
$$\{!\} \cup (\mA'\times S) \to \mA,$$
where $|S|=m+1$ and ``!'' is one of the alphabet symbols. That implies that the size of the code is $$\left\lfloor\frac{\lvert \mathcal{A} \rvert-1}{\lvert S \rvert}\right\rfloor,$$ as desired. 

The decoding works as follows:
\begin{enumerate}
    \item The output state of $V_1$ is $!$ and the output state of $V_2$ is $(y,\beta)$. In this case, the answer is $(y,\ldots,y)$;
    \item The output state of $V_1$ is $(x,\alpha)$ and the output state of $V_2$ is $!$. In this case, the answer is $(x,\ldots,x)$;
    \item The output state of $V_1$ is $(x,\alpha)$ and the output state of $V_2$ is $(y,\beta)$. Here, we check whether $\alpha \ge \beta$. If so, we decode it to $(x,\ldots,x)$. Otherwise, we decode to $(y,\ldots,y)$.
\end{enumerate} 
For the analysis, assume that the $a$ incoming edges of $V_1$ and the $b$ incoming edges of $V_2$ are attacked, so that $a+b\le 2m+1$.
If $a\ge m+1$, then $V_2$ outputs $(y,\beta)$ for some $\beta > m+1$, and thus $y$ coincides with the coordinate of the originally sent codeword. So, if $V_1$ outputs $!$, then the decoding is correct. If $V_1$ outputs $(x,\alpha)$, then we can easily check $\alpha\le m+1 <\beta$ and the decoding is also correct.

Similarly, if $b\ge m+1$, then $V_1$ outputs $(x,\alpha)$ for some $\alpha \ge m+1$, and $x$ coincides with the coordinate of the originally sent codeword. So, if~$V_2$ outputs $!$, then the decoding is correct. If~$V_2$ outputs $(y,\beta)$, then one can easily check $\beta\le m+1 \le \alpha$ and the decoding is also correct. Note that we cannot have  
both $V_1$ and $V_2$ output $!$.

Finally, if $a \le m$ and $b<m+1$, then $V_1$ outputs $(x,\alpha)$ and $V_2$ outputs $(y,\beta)$, where we have~$\alpha,\beta>m+1$. In this case, we have $x=y$, which is the original sent codeword, so that the choice of the decoder does not matter.
\end{proof}

\end{document}
